\begin{tabular}{@{}>{\raggedright\arraybackslash}p{0.16\linewidth}>{\raggedright\arraybackslash}p{0.28\linewidth}>{\raggedright\arraybackslash}p{0.48\linewidth}@{}}
\toprule
Case & Question (start) & Observed \\
\midrule
Disagreement, MATH test:2460 (RWC right, SC wrong); gold \texttt{154} & At a competition with \$N\$ players, the number of players given elite status is equal to \$2\textasciicircum{}\{1+\textbackslash{}lfloor \textbackslash{}log\_\{2\} (N-1) \textbackslash{}rfloor\}-N\$. Suppose that \$19\$ pl\ldots & SC: \texttt{58} (11/32, wrong); Self-certainty + RWC: \texttt{154} (11/32, right). Groups: \texttt{154} 11$\checkmark$; \texttt{58} 11; \textit{no answer} 3; \texttt{45} 1 \\
Disagreement, MATH test:2198 (SC right, RWC wrong); gold \texttt{11} & Let \$g(x) = x\textasciicircum{}2 - 11x + 30,\$ and let \$f(x)\$ be a polynomial such that \textbackslash{}[g(f(x)) = x\textasciicircum{}4 - 14x\textasciicircum{}3 + 62x\textasciicircum{}2 - 91x + 42.\textbackslash{}]Find the sum of all possible values\ldots & SC: \texttt{11} (5/32, right); Self-certainty + RWC: \texttt{14} (4/32, wrong). Groups: \texttt{11} 5$\checkmark$; \texttt{14} 4; \texttt{10\textasciicircum{}\{200\} - 7 \textbackslash{}cdot 10\textasciicircum{}\{1\ldots} 3; \texttt{10\textasciicircum{}\{200\} - 7 \textbackslash{}cdot 10\textasciicircum{}\{1\ldots} 2 \\
Disagreement, MathVista 552 (SC right, RWC wrong); gold \texttt{4} & How many queries have a p-value lower than 0.50? & SC: \texttt{4} (14/32, right); VisConf + RWC: \texttt{2} (9/32, wrong). Groups: \texttt{4} 14$\checkmark$; \texttt{3} 9; \texttt{2} 9 \\
Disagreement, MathVista 643 (RWC right, SC wrong); gold \texttt{A} & As shown in the figure, CD is the diameter of $\odot$O, chord DE $\parallel$ OA, if the degree of $\angle$D is 50.0, then the degree of $\angle$C is () A. 25$^\circ$ B. 30$^\circ$ C. 40$^\circ$ D. 50$^\circ$ & SC: \texttt{C} (11/32, wrong); VisConf + RWC: \texttt{A} (10/32, right). Groups: \texttt{C} 11; \texttt{A} 10$\checkmark$; \textit{no answer} 8; \texttt{D} 3 \\ \midrule
Revision, live test:3504 (wrong to correct); gold \texttt{216} & What is the least positive multiple of 72 that has exactly 16 positive factors? & Initial: \textit{no answer} (1024 tok, truncated); revised: \texttt{216} (566 tok). Revision opens: ``Let's carefully re-examine the steps to ensure we have the correct least positive multiple of 72 that has exactly 16 positive factors. We have already determined that the factorizations of 16 that are relevant for our pr\ldots'' \\
Revision, live test:2892 (changed, still wrong); gold \texttt{-71} & There is a unique polynomial \$P(x)\$ of degree \$8\$ with rational coefficients and leading coefficient \$1,\$ which has the number \textbackslash{}[\textbackslash{}sqrt\{2\} + \textbackslash{}sqrt\{3\} +\ldots & Initial: \textit{no answer} (1024 tok, truncated); revised: \texttt{1} (1332 tok). Revision opens: ``Let's re-evaluate the steps to find the polynomial \textbackslash{}( P(x) \textbackslash{}) with rational coefficients that has \textbackslash{}( \textbackslash{}alpha = \textbackslash{}sqrt\{2\} + \textbackslash{}sqrt\{3\} + \textbackslash{}sqrt\{5\} \textbackslash{}) as a root. 1. **Square \textbackslash{}( \textbackslash{}alpha \textbackslash{}):** \textbackslash{}[ \textbackslash{}alpha\textasciicircum{}2 = (\textbackslash{}sqrt\{2\} + \textbackslash{}sqrt\{3\} + \ldots'' \\
Revision, live test:943 (answer kept); gold \texttt{14} & The arithmetic mean (or average) of \$A\$, \$B\$ and \$C\$ is 10. The value of \$A\$ is six less than the value of \$B\$, and the value of \$C\$ is three more tha\ldots & Initial: \texttt{14} (399 tok); revised: \texttt{14} (311 tok). Revision opens: ``Let's carefully review the solution step-by-step to ensure accuracy. We start with the given information: 1. The arithmetic mean of \textbackslash{}( A \textbackslash{}), \textbackslash{}( B \textbackslash{}), and \textbackslash{}( C \textbackslash{}) is 10. 2. \textbackslash{}( A = B - 6 \textbackslash{}) 3. \textbackslash{}( C = B + 3 \textbackslash{}) From the arit\ldots'' \\ \midrule
Premature stop, MATH test:1803 (agreement); gold \texttt{\textbackslash{}frac35} & In right triangle \$XYZ\$, shown below, what is \$\textbackslash{}sin\{X\}\$? [asy] draw((0,0)--(10,0)--(3.6,4.8)--cycle,black+linewidth(1)); draw(rightanglemark((0,0),(3.\ldots & Stopped at $n=8$ with \texttt{\textbackslash{}frac\{4\}\{5\}} (75\% of revealed, wrong); fixed 32: \texttt{\textbackslash{}frac\{3\}\{5\}} (right). Correct share of full pool 50\%. Groups: \texttt{\textbackslash{}frac\{3\}\{5\}} 16$\checkmark$; \texttt{\textbackslash{}frac\{4\}\{5\}} 15; \texttt{\textbackslash{}frac\{5\}\{4\}} 1 \\
Premature stop, MATH test:1803 (quality); gold \texttt{\textbackslash{}frac35} & In right triangle \$XYZ\$, shown below, what is \$\textbackslash{}sin\{X\}\$? [asy] draw((0,0)--(10,0)--(3.6,4.8)--cycle,black+linewidth(1)); draw(rightanglemark((0,0),(3.\ldots & Stopped at $n=8$ with \texttt{\textbackslash{}frac\{4\}\{5\}} (75\% of revealed, wrong); fixed 32: \texttt{\textbackslash{}frac\{3\}\{5\}} (right). Correct share of full pool 50\%. Groups: \texttt{\textbackslash{}frac\{3\}\{5\}} 16$\checkmark$; \texttt{\textbackslash{}frac\{4\}\{5\}} 15; \texttt{\textbackslash{}frac\{5\}\{4\}} 1 \\
Premature stop, MathVista 434 (agreement); gold \texttt{A} & In the figure above, side AC of triangle ABC is on line l. What is x in terms of k? A. 60-k B. k C. 60+k D. 120-k E. 120-2*k & Stopped at $n=4$ with \texttt{C} (50\% of revealed, wrong); fixed 32: \texttt{A} (right). Correct share of full pool 53\%. Groups: \texttt{A} 17$\checkmark$; \texttt{C} 7; \texttt{B} 4; \textit{no answer} 2 \\
Premature stop, MathVista 434 (quality); gold \texttt{A} & In the figure above, side AC of triangle ABC is on line l. What is x in terms of k? A. 60-k B. k C. 60+k D. 120-k E. 120-2*k & Stopped at $n=4$ with \texttt{C} (50\% of revealed, wrong); fixed 32: \texttt{A} (right). Correct share of full pool 53\%. Groups: \texttt{A} 17$\checkmark$; \texttt{C} 7; \texttt{B} 4; \textit{no answer} 2 \\
\bottomrule
\end{tabular}